% =============================================================================
%  LIVRABLE 1 - Reformulation du sujet sous forme de carte conceptuelle
%  Projet de recherche 4A S7 - Optimisation de modeles neuronaux a oscillateurs couples
%  Compilation : tectonic schema_epure.tex   (ou pdflatex / lualatex)
% =============================================================================
\documentclass[tikz,border=14pt,10pt]{standalone}
\usepackage[utf8]{inputenc}
\usepackage[T1]{fontenc}
\usepackage[french]{babel}
\usepackage{lmodern}
\usepackage{amsmath,amssymb}
\usepackage{enumitem}
\usepackage{tikz}
\usetikzlibrary{arrows.meta,positioning,calc,fit,backgrounds}

% --- PALETTE SOBRE : 1 encre + 4 accents desatures -------------------------
\definecolor{cInk}    {RGB}{26, 32, 44}    % texte
\definecolor{cNavy}   {RGB}{19, 49, 80}    % structure / centre
\definecolor{cMute}   {RGB}{113, 124, 138} % texte secondaire
\definecolor{cLine}   {RGB}{214, 221, 229} % filets & bordures
\definecolor{cOne}    {RGB}{31, 92, 139}   % 1. Modeliser
\definecolor{cTwo}    {RGB}{181, 84, 39}   % 2. Accelerer
\definecolor{cThree}  {RGB}{34, 116, 80}   % 3. Mesurer
\definecolor{cFour}   {RGB}{104, 72, 140}  % 4. Generaliser
\definecolor{cSoft}   {RGB}{248, 250, 252} % fond tres clair

\newcommand{\lead}[1]{\textbf{#1}}
% Tirets et separateurs compatibles avec l'encodage EC (lmodern sans)
\newcommand{\md}{\ \textendash\ }          % tiret de definition
\newcommand{\sep}{\;\textperiodcentered\;} % separateur d'enumeration en ligne
\newcommand{\tdash}{\;\textendash\;}       % tiret de titre

% Puce sobre : tiret court colore
\newlist{dims}{itemize}{1}
\setlist[dims]{leftmargin=2.6mm,labelsep=1.6mm,itemsep=3.0pt,topsep=0pt,parsep=0pt,
               label=\textcolor{cLine!40!cMute}{\rule[0.62ex]{1.5mm}{0.5pt}}}

\begin{document}
\begin{tikzpicture}[
    font=\sffamily,
    >={Stealth[length=2.6mm,width=2mm]},
    % --- conteneur d'une dimension
    dim/.style={
        rectangle, rounded corners=4pt, fill=white, draw=cLine, line width=0.9pt,
        inner sep=10pt, align=left, text=cInk
    },
    % --- pastille numerotee
    tag/.style={
        circle, draw=none, fill=#1, text=white, inner sep=0pt,
        minimum size=6.4mm, font=\sffamily\bfseries\footnotesize,
        draw=white, line width=1.1pt
    },
    % --- anneau d'iteration
    ring/.style={
        draw=cNavy!45, line width=0.9pt, rounded corners=9pt, ->
    },
    ringDash/.style={
        draw=cNavy!45, line width=0.9pt, rounded corners=9pt, ->, dash pattern=on 2.4pt off 2.2pt
    },
    ringLab/.style={
        font=\sffamily\fontsize{7.8}{9.4}\selectfont, text=cNavy!72, inner sep=0pt
    }
]

% =========================================================================
%  TITRE
% =========================================================================
\node (title) [align=center] at (0, 9.75) {
    {\fontsize{15}{18}\selectfont\bfseries\color{cNavy}
     Optimisation de modèles neuronaux à oscillateurs couplés}\\[3pt]
    {\fontsize{9}{11}\selectfont\color{cMute}
     \textsc{Livrable 1}\tdash Reformulation du sujet sous forme de carte conceptuelle}
};

% =========================================================================
%  NOYAU : L'OBJET D'ÉTUDE  (pas de problématique dans le schéma)
% =========================================================================
\node (core) [
    rectangle, rounded corners=5pt, fill=cNavy, draw=none,
    minimum width=6.1cm, inner sep=11pt, align=center, text=white
] at (0, 0.35) {
    \begin{minipage}{5.2cm}
        \centering
        {\fontsize{7.5}{9}\selectfont\bfseries\color{cLine} OBJET D'ÉTUDE}\\[5pt]
        {\fontsize{9.6}{12.6}\selectfont
         Un réseau de \textbf{neurones biophysiques couplés},\\
         simulé \textbf{fidèlement} et calculé \textbf{vite},\\
         jusqu'à grande échelle}
    \end{minipage}
};

% Tension structurante, juste sous le noyau
\node (tension) [
    rectangle, rounded corners=3pt, fill=cSoft, draw=cLine, line width=0.8pt,
    inner xsep=8pt, inner ysep=5pt, align=center, anchor=north
] at ([yshift=-5pt]core.south) {
    {\fontsize{8.4}{10}\selectfont\color{cNavy}%
     \textbf{Vitesse}\;\;$\rightleftharpoons$\;\;\textbf{Fidélité biophysique}}
};

% =========================================================================
%  LES 4 DIMENSIONS
% =========================================================================
% ---------- 1. MODÉLISER (haut gauche) ----------
\node (d1) [dim] at (-7.15, 4.53) {
    \begin{minipage}[t][5.95cm][t]{5.75cm}
        {\fontsize{12}{14}\selectfont\bfseries\color{cOne}MODÉLISER}%
        \hfill{\fontsize{8}{9}\selectfont\color{cMute}\textsc{quoi ?}}\\[1.5pt]
        {\fontsize{8}{9.5}\selectfont\color{cMute}Le substrat biophysique à reproduire}\\[4pt]
        {\color{cOne}\rule{5.75cm}{1.1pt}}\\[6pt]
        {\fontsize{8.2}{10.2}\selectfont\raggedright
        \begin{dims}
            \item \lead{Neurone}\md Hodgkin-Huxley (1952) : état $(V,m,h,n)$,
                  4 EDO raides couplées, pas $dt \le 0,\!05$\,ms
            \item \lead{Synapse}\md modèle cinétique (Destexhe, 1994) :
                  $I_{\mathrm{syn}} = g\,s_{ij}(V-E)$, excitatrice ou inhibitrice
            \item \lead{Connectome}\md graphe Watts-Strogatz (petit monde),
                  matrice creuse $W$, $N : 10 \rightarrow 10^{4}$
            \item \lead{Régimes}\md potentiels d'action, bouffées,
                  synchronisation : la signature à préserver
        \end{dims}}
    \end{minipage}
};

% ---------- 2. ACCÉLÉRER (haut droite) ----------
\node (d2) [dim] at (7.15, 4.53) {
    \begin{minipage}[t][5.95cm][t]{5.75cm}
        {\fontsize{12}{14}\selectfont\bfseries\color{cTwo}ACCÉLÉRER}%
        \hfill{\fontsize{8}{9}\selectfont\color{cMute}\textsc{comment ?}}\\[1.5pt]
        {\fontsize{8}{9.5}\selectfont\color{cMute}Les stratégies de calcul à confronter}\\[4pt]
        {\color{cTwo}\rule{5.75cm}{1.1pt}}\\[6pt]
        {\fontsize{8.2}{10.2}\selectfont\raggedright
        \begin{dims}
            \item \lead{Référence CPU}\md intégration exacte à pas adaptatif
                  (SciPy \texttt{solve\_ivp}) : la vérité terrain
            \item \lead{Portage GPU}\md parallélisme \textsc{simt}, un fil par
                  neurone : PyTorch, CuPy, noyaux CUDA
            \item \lead{Leviers}\md vectorisation, produit creux (SpMV),
                  pas fixe, occupation et transferts VRAM
            \item \lead{Compromis}\md précision numérique
                  (\texttt{float32} / \texttt{float64}) contre débit
        \end{dims}}
    \end{minipage}
};

% ---------- 3. MESURER (bas droite) ----------
\node (d3) [dim] at (7.15, -3.83) {
    \begin{minipage}[t][5.95cm][t]{5.75cm}
        {\fontsize{12}{14}\selectfont\bfseries\color{cThree}MESURER}%
        \hfill{\fontsize{8}{9}\selectfont\color{cMute}\textsc{combien ?}}\\[1.5pt]
        {\fontsize{8}{9.5}\selectfont\color{cMute}La preuve expérimentale à produire}\\[4pt]
        {\color{cThree}\rule{5.75cm}{1.1pt}}\\[6pt]
        {\fontsize{8.2}{10.2}\selectfont\raggedright
        \begin{dims}
            \item \lead{Gain}\md accélération
                  $S = T_{\mathrm{CPU}}/T_{\mathrm{GPU}}$ en fonction de $N$
            \item \lead{Passage à l'échelle}\md coût par pas $\Delta t$, mémoire
                  occupée, seuil de bascule CPU\,$\rightarrow$\,GPU
            \item \lead{Non-dérive}\md
                  $\lVert V_{\mathrm{GPU}} - V_{\mathrm{CPU}} \rVert_\infty < \varepsilon$,
                  dates de spikes et synchronisation préservées
            \item \lead{Protocole}\md mêmes conditions initiales, exécutions
                  répétées, matériel et versions documentés
        \end{dims}}
    \end{minipage}
};

% ---------- 4. GÉNÉRALISER (bas gauche) ----------
\node (d4) [dim] at (-7.15, -3.83) {
    \begin{minipage}[t][5.95cm][t]{5.75cm}
        {\fontsize{12}{14}\selectfont\bfseries\color{cFour}GÉNÉRALISER}%
        \hfill{\fontsize{8}{9}\selectfont\color{cMute}\textsc{pour quoi ?}}\\[1.5pt]
        {\fontsize{8}{9.5}\selectfont\color{cMute}Le socle réutilisable à livrer}\\[4pt]
        {\color{cFour}\rule{5.75cm}{1.1pt}}\\[6pt]
        {\fontsize{8.2}{10.2}\selectfont\raggedright
        \begin{dims}
            \item \lead{Briques interchangeables}\md \textsf{[neurone]}
                  \textsf{[synapse]} \textsf{[connectome]} \textsf{[solveur]}
            \item \lead{Généricité}\md changer de modèle ou de topologie
                  sans réécrire le moteur de simulation
            \item \lead{Livraison}\md bibliothèque documentée, tests de
                  non-régression, notebooks de démonstration
            \item \lead{Débouchés}\md neurosciences computationnelles,
                  bio-électronique, interfaces cerveau-machine
        \end{dims}}
    \end{minipage}
};

% ---------- pastilles numérotées ----------
\node [tag=cOne]   at ([shift={(0mm,0mm)}]d1.north west) {1};
\node [tag=cTwo]   at ([shift={(0mm,0mm)}]d2.north west) {2};
\node [tag=cThree] at ([shift={(0mm,0mm)}]d3.north west) {3};
\node [tag=cFour]  at ([shift={(0mm,0mm)}]d4.north west) {4};

% =========================================================================
%  ANNEAU D'ITÉRATION  1 -> 2 -> 3 -> 4 -> 1
% =========================================================================
\def\ringT{8.70}   % ordonnée du segment haut
\def\ringB{-8.00}  % ordonnée du segment bas
\def\ringR{11.15}  % abscisse du segment droit
\def\ringL{-11.15} % abscisse du segment gauche

\draw [ring] (d1.north) -- (-7.15,\ringT) -- (7.15,\ringT) -- (d2.north);
\draw [ring] (d2.east)  -- (\ringR,4.53) -- (\ringR,-3.83) -- (d3.east);
\draw [ring] (d3.south) -- (7.15,\ringB) -- (-7.15,\ringB) -- (d4.south);
\draw [ringDash] (d4.west) -- (\ringL,-3.83) -- (\ringL,4.53) -- (d1.west);

% Libelles de liaison : tous horizontaux, poses le long du trait,
% a l'interieur de l'anneau, sans cartouche.
\node [ringLab, anchor=north]                        at (     0,  8.50) {équations, graphe et jeu de tests};
\node [ringLab, anchor=south]                        at (     0, -7.80) {temps mesurés et traces $V(t)$};
\node [ringLab, anchor=east, align=right, text width=3.5cm]
      at ( 10.90, 0.35) {exécutions comparées,\\ CPU contre GPU};
\node [ringLab, anchor=west, align=left,  text width=3.5cm]
      at (-10.90, 0.35) {itérer : autre modèle,\\ autre topologie};

\end{tikzpicture}
\end{document}
